% V24 §16 early independent collation report; not author text.
\begingroup
\def\Sym{\operatorname{Sym}}
\def\gr{\operatorname{gr}}
\providecommand{\Oa}{\mathcal O}
\providecommand{\Pp}{\mathcal P}
\providecommand{\Tt}{\mathcal T}
\head{第二十四版专项核查：第二篇第16节前半}
\noindent 本报告对应法稿物理第95--103页、合并英文稿第62--68页至可逆模图，以及独立\texttt{EGAV1-2.ps}全部6页的重合文字。以下“补证”“反例”“类型核对”均为本版译校工作，不是原作者已经写出的证明。法稿的有效正文、删除文字、后续术语修订和英文录入错误分别处理；不由编译通过推论数学正确。

\audit{V24-EAR-A01}{来源、边界与错误归属}{正文接续和三载体边界已逐页对读}
本模块从第二篇标题、定义1开始，止于可逆模$K$的图（$D_K$）。FR104、EN68图后的英译划去说明及增广定义属于后半模块；命题8不属于本半。PS实际转换后为6页：标题页印刷页码1，随后正文印刷页码1--5；它止于行列式交换图前的“交换性”一句，没有图本身及后续内容。本版并未把这个载体截断误认为法文缺文。

\begin{longtable}{@{}p{.18\textwidth}p{.33\textwidth}p{.39\textwidth}@{}}
\toprule
位置 & 实际证据 & 本版处理\\\midrule
\endfirsthead
\toprule 位置 & 实际证据 & 本版处理\\\midrule\endhead
FR95定义1 & “非退化”的“非”已被作者划去；$\Sym$漏二次上标 & 采用“退化”；补$\Sym^2$；不归为未改的英译错误\\
FR95注2后 & 可辨的一整块重数断言已划去，英文未载 & 在EAR006中完整译出可辨数学内容并给反例，保留删除地位\\
FR95定义3 & FR、EN62及PS均把域值点写成$k$值 & 改为$k'$值；注明法稿已有笔误\\
FR96--97 & 命题5及二阶初始项的相对底错写$S$ & 恢复$Y$，与尚未引入$S$的上下文一致\\
FR98维数 & 法稿正确为$(m+d)-(d+1)$；EN64、PS为$m+1$ & 确认为英文载体的录入错误\\
FR98图 & FR及PS的$\Omega_{X/S}^1$、$\Omega_{Y/S}^1$正确；EN64错录 & 依FR、PS恢复完整图，不统一抹去载体差异\\
FR99 & 层单射/纤维秩混淆、相对底和余法列末项有误 & 正文给正确Jacobian条件及$\Omega_{W/S}^1$，原句与反例入脚注\\
FR99主部分解 & 法稿及英文都误用$X/S$ & 根据图中对象及秩恢复$X/Y$；切层英编疑问保留\\
FR100--101 & 法稿外幂、$p+q=m+1$、$M/Q$可辨正确；英文多处错录 & 恢复目标秩外幂及正确商模；明确局部直和因子条件\\
FR101 & 图中$(\det M)^{-1}$清楚；正文上标难辨 & 按图及类型确定逆线丛，不强认难辨上标\\
FR102 & 第二正合列末项秩错误来自法稿 & 正文修为$\Omega_{X/Y}^1|_U\oplus\Omega_{V'/S}^1|_U$\\
FR102--103 & 正规性的阈值、非闭点孤立性缺条件来自法稿 & 分别补超曲面维数至少2和闭点限制，并给反例\\
FR103图 & EN写$K^{-1}$；法稿指数覆写不清 & 用行列式核对改为$K^{-(d+1)}$，原字仍列未定\\
\bottomrule
\end{longtable}

上述“FR”一律是PDF物理页序。实际观看记录另存于本版\texttt{SOURCE\_MAP.md}和源页核验记录，不将文本抽取、历史截图或图片文件存在当作已经目视。

\audit{V24-EAR-A02}{二次初始项、极化、重数与零秩边界}{两种二次型性质必须分开；删除段的重数等价不成立}
设$A$为正则局部环，极大理想$\mathfrak m$，余域$\kappa$。有
\[
 \gr_{\mathfrak m}A\simeq\Sym_{\kappa}(\mathfrak m/\mathfrak m^2).
\]
因此$f\in\mathfrak m^2$的二次初始项$q$属于$\Sym^2(\mathfrak m/\mathfrak m^2)$；它给出对偶切空间上的二次型。其极化为
\[
 b_q(v,w)=q(v+w)-q(v)-q(w).
\]
本文本使用$T\to T^*$的$b_q$是否同构判定非退化。整个定义无须$2$可逆。

特征2时，取$q=X^2+YZ$，按基$X,Y,Z$，极化矩阵为
\[
 H_q=\begin{pmatrix}0&0&0\\0&0&1\\0&1&0\end{pmatrix}.
\]
它的根基为$X$方向，故退化。然而射影曲线$q=0$的梯度是$(0,Z,Y)$；梯度全零将迫使$Y=Z=0$，方程再迫使$X=0$，这在射影空间中不是一点。因此射影二次曲线光滑。这一直接计算说明作者FR112--113所说的术语冲突，正文遵从其后续改名，但没有把“极化非退化”替换成“光滑二次型”。后续作者术语说明本轮只用于核证，不计作本半新增全文翻译。

FR95删除段的“重数2当且仅当普通奇异”也不能成立。例如$k$特征0，
\[
 B=k[x,y]_{(x,y)}/(y^2-x^3).
\]
初始关系为$y^2$，所以$\gr_{\mathfrak m}B=k[X,Y]/(Y^2)$，其Hilbert--Samuel重数为2；极化却在$X$方向有根基。重数2只能约束最低非零次数，不能强制二次型非退化。已划去段不是本版维持的有效断言。

另有不可省略的零秩边界：若$A=\kappa$、$f=0$，则切空间为零，$0\to0$是同构，$\det$空矩阵为1。按定义1字面，它是奇异而非超奇异零点；但$V(f)=\Spec\kappa$本身正则，且没有维数下降1。这不允许从定义1无条件推出“非零二次超曲面奇点”。本版保留定义的字面代数约定，在开始把它解释为实际超曲面奇点时明确要求正余切秩，并与后半非退化二次型增广的正秩条件一致。

\audit{V24-EAR-A03}{一阶主部、常数项与余法映射}{常数因子不可丢失；光滑性取决于余域纤维满秩}
\noindent\textbf{译校补证。}一阶主部有典范分解
\[
 \Pp^1_{X/Y}\simeq\Oa_X\oplus\Omega^1_{X/Y},
 \qquad d^1\varphi=(\varphi,d_{X/Y}\varphi).
\]
因此$V'=V(d^1\varphi)$由函数$\varphi$及相对一阶导数共同定义。只取导数的零点集会丢去“零点”条件；例如$\varphi=1$时相对导数为零，但$V'$为空。

局部选取$X/Y$的相对étale坐标$t_1,\ldots,t_d$，写$g_i=\partial_i\varphi$。在图（D）的余法映射
\[
 A=(\Pp^1_{X/Y})^\vee|_{V'}
 \longrightarrow B=\Omega^1_{X/S}|_{V'}
\]
中，常数因子$\Oa_{V'}$的生成元映为$d_{X/S}\varphi$，而切向因子第$i$个生成元映为$d_{X/S}g_i$。映到$C=\Omega^1_{X/Y}|_{V'}$后，第一项因$g_i=0$而为零；余下的矩阵为
\[
 H=(\partial_j\partial_i\varphi)_{i,j}.
\]
这就是二次初始项的极化矩阵：平方项$cT_i^2$贡献对角元$2c$，混合项$cT_iT_j$贡献两个非对角元$c$。没有除以2。

若$V'/S$光滑且有预期余维$d+1$，余法映射$A\to B$在每个余域纤维上单射。因此其常数因子生成元的像$d_{X/S}\varphi$在每个纤维上非零。这才证明原NB所说的“相对于$Y$奇异的零点，相对于$S$非奇异”，而不是因为$\mu$杀掉常数因子就把它的环境微分也设为零。图中的$\Oa_{V'}\to Q$正是这个微分落入$Q=\Omega^1_{Y/S}|_{V'}$后的因子化。

对于一般局部自由模$F$的截面$\psi$，设$J$是它的零点概形$W$的理想。先有满射
$F^\vee|_W\twoheadrightarrow J/J^2$，再与余法映射复合，得到正文的右正合列。若$X/S$光滑，$F$秩$r$，在$w$处的Jacobian矩阵余域秩等于$r$，等价于$W/S$光滑且相对维数比环境少$r$。反之，只知层同态单射不足以推出这些结论。

具体地，在上述尖点整环$B$上，
\[
 B\longrightarrow B^2,\qquad h\longmapsto h(-3x^2,2y)
\]
单射，因为$B$为整环且向量有非零分量；但在原点余域纤维上这是零映射。原文“能够在左端添零”的措辞不能混同这两种单射性。

\audit{V24-EAR-A04}{分歧子概形与行列式线丛}{所核的是理想相等，不只是底层集合相等}
图（D）给出局部自由模中的两条局部直和包含$P,Q\subset M$。在几何应用中，
\[
 \rk M=d+m,\quad\rk P=d+1,\quad\rk Q=m,
\]
所以$\rk P+\rk Q=\rk M+1$。设一般秩为$m_0,p,q$，并令$S_0=M/Q$、$R=M/P$。商模局部自由是以下行列式计算的必要条件。

映射$P\to S_0$的目标秩为$p-1$；其最高子式给出
\[
 a\in P\otimes\det S_0\otimes(\det P)^{-1}
    =P\otimes L^{-1},
\]
其中
\[
 L=\det P\otimes\det Q\otimes(\det M)^{-1}.
\]
同样得到$b\in Q\otimes L^{-1}$。选择相容的行列式符号约定后，$a,b$在$M\otimes L^{-1}$中的像相等；不固定约定时至多相差$-1$。其局部内容可用余子式展开验证：把$P\to M/Q$写成$(p-1)\times p$矩阵，带交替符号的最大子式组成核向量；包含到$M$后它落入$Q$，并正是另一个商映射的对应余子式向量。

因为$P,Q$的包含均有局部左逆，把$a$视为$P$中的向量或$M$中的向量，其坐标理想相同；$b$亦然。这给出$\mathscr I(a)=\mathscr I(b)$。在共同零点之外，余子式向量有一个单位分量，故$L$的像为一条直和因子，并等于$P\cap Q$。

应用于$P=(\Pp^1_{X/Y})^\vee|_{V'}$时，$P\to M/Q$的矩阵是$(H\;0)$；最大子式除$\det H$外都为零。因此其理想为$(D)$。另一方面，$Q\to R=\Omega^1_{V'/S}$的最大子式定义源端分歧子概形，亦即
\[
 \operatorname{Fitt}_0(\Omega^1_{V'/Y}).
\]
上述理想等式便严格给出
\[
 V(\varphi)^{\mathrm{sup\,sing}}
   =\operatorname{Ram}(V'/Y)
\]
作为闭子预概形的相等。该论证在特征2仍成立，因为始终使用极化矩阵，并未改成射影二次曲面光滑判据。

\audit{V24-EAR-A05}{第二条正合列的修正及共同线}{旧附记的条件性解释已保留并落实到正文}
令$U=V'\setminus V(\varphi)^{\mathrm{sup\,sing}}$，在$U$上写
\[
 E=\Omega^1_{X/S}|_U,\quad
 P=(\Pp^1_{X/Y})^\vee|_U,\quad
 Q=\Omega^1_{Y/S}|_U.
\]
由上一项，$P+Q=E$，共同线$K_0=P\cap Q\simeq\Oa_U$。自然映射
\[
 E\longrightarrow E/Q\oplus E/P
\]
是满射，核为$K_0$。满射性可直接验证：给定代表元$e_1,e_2$，将$e_1-e_2$写成$p+q$，其中$p\in P$、$q\in Q$，则$e=e_1-q=e_2+p$同时代表所需两个商类。

也可沿旧附记原有的构造，取
\[
 (P/K_0)\oplus(Q/K_0)\longrightarrow E/K_0,
 \qquad(\bar p,\bar q)\longmapsto\overline{p+q}.
\]
和条件给满射；若$p+q\in K_0$，则$p\in Q\cap P=K_0$，随后$q\in K_0$，故核为零。再由
\[
 P/K_0\simeq E/Q\simeq\Omega^1_{X/Y}|_U,
 \qquad Q/K_0\simeq E/P\simeq\Omega^1_{V'/S}|_U
\]
即得正文修正后的正合列。右侧秩为$d+(m-1)$，与中项秩$d+m$及核秩1一致。

原法稿直接写$\Omega^1_{Y/S}|_U$作为第二个因子，少商掉这条共同线。前版已把此问题写入附记，但仍将错误式留在正文；按用户已经确定的“正文写正确、原错误留脚注”的要求，本版把有明确条件和证明支持的订正式落实到正文，同时保存原式、秩矛盾及这一补证。

\audit{V24-EAR-A06}{单位倍与可逆模主部图}{完整扭曲为$K^{-(d+1)}$，原手写指数不强辨}
局部平凡化可逆模后，由
\[
 d(u\varphi)=u\,d\varphi+\varphi\,du
\]
可知，限制到$V(\varphi)$后，一阶微分只乘以$u$。再限制到$V'$，$\varphi$及相对一阶导数均为零，二阶初始项中的其他乘积项也消失，故
\[
 M(u\varphi)=uM(\varphi),\qquad
 D(u\varphi)=u^dD(\varphi).
\]
这些等式是所列限制上的等式，不是整个$X$上二阶导数的无条件等式。单位因子保证相应零理想不变。

对可逆模$K$，不作典范拆分，而使用主部正合列
\[
 0\longrightarrow\Omega^1_{X/Y}\otimes K
 \longrightarrow\Pp^1_{X/Y}(K)\longrightarrow K\longrightarrow0.
\]
由此
\[
 \det(\Pp^1_{X/Y}(K))
       =\omega_{X/Y}\otimes K^{d+1}.
\]
取对偶，再与$\det Q$及$(\det B)^{-1}$相乘，得到
\[
 L_K=(\omega_{X/Y}^{-2}\otimes K^{-(d+1)})|_{V'}.
\]
相对二阶初始项属于$\Sym^2\Omega^1_{X/Y}\otimes K$，其极化可视为
\[
 \Tt_{X/Y}\otimes K^{-1}\longrightarrow\Omega^1_{X/Y}.
\]
其行列式属于$\omega_{X/Y}^{\otimes2}\otimes K^d$；另一方面
\[
 \mathscr Hom(L_K,K^{-1}|_{V'})
  =(\omega_{X/Y}^{\otimes2}\otimes K^d)|_{V'},
\]
正与左下箭头类型一致。它也与单位倍$u^d$吻合。

FR103的实际4倍裁切仍可见指数覆写，不能断言逐字读出了完整指数。这里的订正依据是明确的主部正合列、行列式关系和单位缩放，不是对模糊笔迹的猜测。原图的箭头方向、顶端和底端零对象全部保留。

\audit{V24-EAR-A07}{孤立性与正规性：两个不同漏洞}{修正维数阈值；区分局部严格泛化与全局拓扑孤立}
\noindent\textbf{域上光滑环境的独立补证。}
设$X/k$光滑，在所论点的相对维数$d>0$，并取几何奇异而非超奇异零点$x$。在相对étale坐标中，令$g_i=\partial_i\varphi$。在$x$处，$dg_i$相对于$dt_j$的矩阵就是可逆Hessian$H$。Jacobian判据说明$V(g_1,\ldots,g_d)$在$x$附近对$k$étale，因此零维，且其点的余域对$k$有限可分。因为$\varphi(x)=0$，在这个局部零维étale环中$\varphi$也为零，故$x$附近的临界零概形$V(\varphi,g_1,\ldots,g_d)$就是这个étale零点概形。这证明该点在非光滑点集合中孤立，也解释此处与一般非闭点局部定义的区别。

等价地，在几何域值点$x'$处取光滑参数，$g_i$的线性项矩阵为$H$；由Nakayama引理，$g_i$生成该点的极大理想。该补证对任意特征成立，并没有使用$2$的逆。

\noindent\textbf{正规性的正确维数。}
在几何纤维中，非退化二次初始项在$d>0$时非零，故$\varphi$是正则局部环中的非零元素，商环是维数$d-1$的超曲面局部环并为Cohen--Macaulay。若$d-1\ge2$，孤立奇点以外的点正则，所以所有高度至多1的局部化都正则，即满足$R_1$；Cohen--Macaulay给出$S_2$，于是正规。这里用的标准正规性判据可核对\href{https://stacks.math.columbia.edu/tag/031S}{Stacks Project，Tag 031S}；概形版本为\href{https://stacks.math.columbia.edu/tag/033P}{Tag 033P，Lemma 28.12.5}。应用于各几何扩域给出所述几何正规性，而不是只验证原余域上的正规。

原$d\ge2$在原文的光滑塔假设内也有反例，不能用增加原文已经具备的条件救回。取
\[
 S=\Spec k,\quad Y=\mathbb A^1_k,
 \quad X=\mathbb A^2_k\times Y,\quad\varphi=uv-t.
\]
则$X/Y$相对维数$d=2$、$Y/S$维数$m=1$，$V'$是单个光滑点，恰有相对维数$m-1=0$；极化矩阵行列式为$-1$，在任意特征都可逆。但$t=0$纤维在原点是$uv=0$，局部环有两个分支，不是正规局部环。所以阈值必须落在超曲面维数至少2，即环境$d\ge3$。

\noindent\textbf{非闭点的拓扑漏洞。}
法稿另在一般正则预概形语境中宣布：非退化二次零点的每个严格泛化点正则。这是局部环去掉闭点后的断言；不能仅凭奇异集闭，就推断该点在全局奇异集中孤立。反例取
\[
 X=\Spec k[u,v,w],\quad\varphi=uv,\quad
 x=(u,v),\qquad\operatorname{char}k\ne2.
\]
$x$是奇异直线的泛点，$\Oa_{X,x}=k(w)[u,v]_{(u,v)}$，$uv$的极化非退化。严格泛化点只有相应分支的泛点，均正则；但奇异集是闭直线$V(u,v)$，其泛点不是孤立点。$X$甚至是优良概形。只有再加$x$为该闭奇异集的闭点，才能在局部诺特情形由没有严格泛化、也没有严格特化推出孤立。

\noindent\textbf{本版仍保留的证明限度。}
本版已完整补证域上光滑情形及正规性维数修正，并给出上述反例。法稿对任意正则局部环境的严格泛化正则性，只写“可以证明”，没有展开证明。本版不声称通过域上Hessian论证就已经覆盖所有混合特征或非优良正则局部环；这个一般陈述的独立全范围证明仍列待核。此限度不影响两个反例及正文的必要条件订正。

\audit{V24-EAR-A08}{前版附记与补证的保全}{旧文件完整停用保留，新版追踪而非静默删除}
旧\texttt{text/v1\_2.tex}的61,675字节全文应保持原字节存档。其前半正文和本半旧附记的内容对应如下：
\begin{longtable}{@{}p{.37\textwidth}p{.53\textwidth}@{}}
\toprule 旧内容块 & 新正文及报告去向\\\midrule\endfirsthead
\toprule 旧内容块 & 新正文及报告去向\\\midrule\endhead
定义1至注4(c)的V-H-C16及D17-S13注 & EAR003--008、EAR034；A02的特征2、重数及零秩核查\\
命题5、相对底及一二阶初始项 & EAR009--012；A03的常数项及微分关系\\
Jacobian层单射/纤维单射与余法列 & EAR016--019；A03完整保留尖点反例及纤维秩条件\\
秩图、最高子式及行列式线丛 & EAR015、EAR021--025；A04保留理想而非底集的核对\\
“符号与法文核对” & EAR026--027；区分$\omega/W$、$\otimes/\times$、$S/s$与对偶记号\\
“第二条正合列的秩问题” & EAR027；A05保留原式的秩矛盾\\
“按前图得到的条件性解释” & EAR027；A05保留共同线、和条件、商模同构及逐项核证\\
“原文对读与记号” & EAR013、EAR028；未补造命题6/总结命题，保留限制微分含义\\
“单位变换的适用位置” & EAR028；A06完整保留乘积法则与二阶限制条件\\
“措辞及范围” & EAR004、EAR030--031；严格泛化更正、作者术语修订及新闭点反例\\
“正规性的维数疑点” & EAR029、EAR032；A07将已确证订正移入正文，补在原光滑塔假设内的反例\\
“图的对读”及“左下角的指数” & EAR033；A06保留全部箭头与行列式推导，法稿覆写仍未定\\
“划去标记”及旧第56--62块 & 由后半模块接续，具体条目见其继承映射；本半不宣称重审后半证明\\
\bottomrule
\end{longtable}

本次改变的是确证错误在正文与译校注之间的位置，并增加错误归属和必要边界；已有可用补证没有因替换正文输入而丢弃。前版“暂不替换错误正文”的保守说明已被本版的明确订正取代，其历史原字仍在停用旧源中。独立PS载体没有图后内容，故该部分继承以实际法文及合并英文为据，不声称所有内容都有三载体互证。
\endgroup
